\clearpage\SourceLeaf{151}%
\pdfbookmark[1]{Gauß an von Lindenau}{lindenau1808aug}
\fancyhead[L]{\small GAUSS AN VON LINDENAU}
\begin{center}
{\small GAUSS AN VON LINDENAU.\par}
\medskip\hrule\medskip
{\small Monatliche Correspondenz zur Beförderung der Erd- und Himmels-Kunde,\par
herausgegeben vom Freiherrn \textsc{von Zach}. September 1808.\par}
\medskip\hrule
\end{center}\medskip
\begin{flushright}Göttingen, 30. August 1808.\end{flushright}
\begin{center}\small
—\ —\ —\ —\ —\ —\ —\ —\ —\ —\ —\ —\ —\ —\ —\ —\ —
\end{center}
Ich habe die Ehre, Ihnen hier für die \emph{Mon. Corresp.} [S.\ 141 dieses Bandes] einen Aufsatz über eine Aufgabe zu schicken, deren praktische Anwendung sehr zu empfehlen ist. Mir scheint, dass dieses Verfahren zur Bestimmung der Polhöhe in der Praxis eine sehr hohe Genauigkeit zulässt und die zuverlässigsten Ergebnisse liefert, wenn man keine Gelegenheit hat, mit Repetitionskreisen zu beobachten. Zumindest entfallen alle Umstände, die beim üblichen Verfahren mit Sonnenhöhen die Ergebnisse bisweilen fraglich machen können. Das sind das Ablesen, Fehler des Sextanten, Blendgläser, das Glasdach und die Refraktion. Zwar liefert jede Kombination von drei Sternen immer nur ein Ergebnis. Doch wie sehr man die Beobachtungen der Meridianhöhen der Sonne auch vervielfacht, bleiben die aus den erwähnten Umständen herrührenden Fehler mit Ausnahme des ersten immer in voller Stärke bestehen. Außerdem hindert einen nichts daran, dieselben Beobachtungen in mehreren Nächten zu wiederholen und selbst in einer Nacht mehrere Kombinationen zu bilden, wenn man sie so auswählt, dass die drei Beobachtungen in nicht zu großen Abständen aufeinanderfolgen. Letzte Nacht habe ich die drei Sterne, die im Reispiel der Abhandlung angeführt sind, noch einmal in
\SourceLeaf{152}%
\expandafter\def\csname HistoricalNoteText7\endcsname{Fehler des Sextanten bei $105^0\,52'= -49''4$.}%
einer etwas anderen Höhe beobachtet und für die Polhöhe $51^0\,31'\,54''4$\HistoricalNote{7} gefunden. Die drei Ergebnisse lauten:
\begin{center}
\begin{tabular}{@{}r@{\ }r@{\quad}c@{\quad}l@{}}
August & 25 & $.$ & $51^0\quad31'\quad56''7$\\
       & 27 & $.$ & $51\quad31\quad51.5$\\
       & 30 & $.$ & $51\quad31\quad54.4$\\\cline{2-4}
\multicolumn{2}{@{}r@{\quad}}{im Mittel} & $.$ & $51\quad31\quad54.2$
\end{tabular}
\end{center}
genau der Wert, den \textsc{Mayer} für unsere Polhöhe ermittelt hat. Auch wenn man diesen Grad der Übereinstimmung hier zum Teil für zufällig halten muss, glaube ich doch, dass man durch häufigeres Wiederholen und Vervielfachen der Beobachtungen mit gut bestimmten Sternen der wahren Polhöhe immer bis auf wenige Sekunden nahekommen kann. Alles hängt allein von der Vergrößerung und Klarheit des Fernrohrs und der Aufmerksamkeit des Beobachters ab. Aus Meridianhöhen der Sonne, die stets ausgezeichnet miteinander übereinstimmten, hatte ich zur Zeit der Sonnenwende mit einem Sextanten mehrmals $51^0\,32'\,32''$ erhalten, also ein um mehr als eine halbe Minute fehlerhaftes Ergebnis. Ob der Fehler meines Sextanten allein daher rührt, dass der ganze Bogen zu groß ist, habe ich noch nicht hinreichend geprüft. Gerade mit der von mir empfohlenen Methode wird sich das künftig aber sehr leicht prüfen lassen.

\par\bigskip\centerline{\rule{.27\linewidth}{.4pt}}
